%% P20 --- Stochastic optimisation and differentiating through randomness
%%
%% Stub. The brief below is the contract for this program: what it must argue,
%% and what it must buy the reader. Writing the program means deleting the
%% \programstub{} block and replacing it with the program. If the brief turns
%% out to be wrong once the mathematics has been worked, change it and record
%% the contradiction in CLAUDE.md under *Resolved questions*.

\program{Stochastic optimisation and differentiating through randomness}
\label{prog:P20}

\programstub{%
Argument: a minibatch gradient is an unbiased estimator with variance
proportional to \code{1/B}, and the noise is a property of the algorithm
rather than a defect in it. Payoff: why the loss curve is noisy and what a
moving average of it hides; the linear scaling rule for batch size and
learning rate, presented as widely repeated folklore with a limited
empirical basis, not as a law; gradient clipping as a bound on the update
rather than on the gradient; and the two ways to get a gradient through a
sampling step --- the score-function estimator (REINFORCE, unbiased and high
variance) and the reparameterisation trick (low variance, requires a
differentiable path) --- which is the fork that separates policy-gradient
RLHF from a VAE.

\medskip
\textbf{Planned length:} 50 frames.
}
